\documentclass[12pt]{article}
\usepackage[utf8]{inputenc}
\usepackage[T1]{fontenc}
\usepackage{hyperref}
\title{Saweru Language Resources for K-12 Curriculum Development}
\author{Compiled by Claude Sonnet 5.5 (claude-sonnet-5-5) for Andr\'as Kornai}
\date{}
\begin{document}
\maketitle

\section{Introduction}
Saweru (ISO 639-3: swr; Glottocode sawe1240) belongs to the small Yawa-Saweru family of Papua, Indonesia, together with Yawa. It is spoken on Serui Island in Cenderawasih Bay. Wikipedia says it is moribund, with about 150 speakers in an ethnic group of 300; Joshua Project says a few hundred speakers remain. Sources disagree on whether it is a separate language or a Yawa dialect. Younger people increasingly use Indonesian.

\section{Primary Educational Resources}
The searches found no primers, textbooks or school programmes for Saweru.

\section{Linguistic and Reference Materials}
\begin{itemize}
\item Anne van Schie, \emph{The Papuan languages of Yapen Island: history and typology} (Vrije Universiteit Amsterdam). It covers the Yawa-Saweru family, which has been largely undescribed: \url{https://research.vu.nl/en/publications/the-papuan-languages-of-yapen-island-history-and-typology/}
\item LOT publication notice for the same work: \url{https://lotschool.nl/?p=12437}
\item LINGUIST List issue 37.2414 (appeared in search results for this family; I did not open it): \url{https://linguistlist.org/issues/37/2414}
\item Wikipedia entry for Saweru: \url{https://en.wikipedia.org/wiki/Saweru_language}
\item Wikipedia entry for the Yawa languages: \url{https://en.wikipedia.org/wiki/Yawa_languages}
\end{itemize}

\section{Religious and Community Resources}
\begin{itemize}
\item Joshua Project profile, Saweru in Indonesia (written from a missionary perspective): \url{https://joshuaproject.net/people_groups/14792}
\end{itemize}

\section{Digital and Audio Resources}
The searches found no audio, video or software resources for Saweru.

\section{Curriculum Development Context}
Saweru is very poorly documented, and no educational materials were found. Documentation with the remaining speakers would have to come first. Because Saweru and Yawa are closely related, any Yawa materials might be relevant, but I did not search for them.
\end{document}